\subsubsection{Buying a claim whose credit loss is already in the price}
\label{sec:bank-poci-rollforward}

The preceding loan became credit-impaired after origination. A purchased or
originated credit-impaired asset, abbreviated POCI, has a different starting
point: expected credit losses already enter the cash flows used to determine
its credit-adjusted effective interest rate. Charging those same initial losses
again would reduce the asset twice.

Suppose the bank buys a fixed claim of 121, due in two annual periods, for
50. Its initial expected collection is 60.5. Assume separately established
POCI status, amortised-cost classification and acquisition fair value equal to
the price, with no fees, interim payments, modifications, foreign exchange or
tax. Assume that reasonable prospects of fuller recovery remain until the
terminal assessment, so no earlier full or partial write-off is required.
Expected cash below the contractual amount does not establish that some
portion has no reasonable prospect of recovery. Paragraphs 5.4.4 and B5.4.9
require that separate assessment. The original credit-adjusted rate \(r\),
with positive compounding factor \(1+r\), must discount expected
collection to the amount paid:
\[
 50=\frac{60.5}{(1+r)^2},\qquad r=10\%.
\]
The contractual claim remains 121. Its discounted expected shortfall is
\((121-60.5)/1.1^2=50\), but this is already incorporated in the
purchase measurement and rate. The additional loss allowance at acquisition
is zero. IFRS 9.5.5.13 recognizes only subsequent cumulative changes in
lifetime expected credit losses for POCI assets; paragraph 5.4.1(a) applies
the credit-adjusted rate to amortised cost from initial recognition. These
are the inspected historical EU 2023/1803 provisions, together with the
Appendix A rate definition and B5.4.7 \citep{eu20231803}. A large discount
or high credit risk alone does not establish POCI status.

With unchanged expectations, first-year interest is \(50\times10\%=5\),
leaving carrying amount 55. Now suppose a supported credit reassessment
raises expected terminal collection to 88. With one year remaining, its
present value at the original rate is \(88/1.1=80\). The increase
\(80-55=25\) is an impairment gain. The following year's interest is
\(80\times10\%=8\), bringing carrying amount to 88 if expectations
remain unchanged. Improved credit does not replace the original rate or
switch this POCI asset to the ordinary gross-interest method. The cure rule
in paragraph 5.4.2 refers specifically to assets that became impaired
\emph{after} initial recognition under paragraph 5.4.1(b).

To reconcile the allowance as well as the net asset, distinguish three
amounts. Let \(G_t\) be carrying amount before the subsequent loss
adjustment, \(L_t\) the signed allowance, and \(N_t=G_t-L_t\) amortised
cost after that adjustment. With no receipts before maturity, the original
purchase measurement accretes to \(G_t=50(1.1)^t\). If \(E_t\) is
expected terminal collection at annual date \(t\), then
\[
 N_t=\frac{E_t}{1.1^{2-t}},\qquad
 L_t=G_t-N_t
     =\frac{(121-E_t)-(121-60.5)}{1.1^{2-t}}.
\]
The last expression compares current and initial expected shortfalls at
the same reporting date. It explains both signs. At year one, expected
collection of 44 gives \(N_1=40\) and \(L_1=15\), an impairment loss
of 15. Expected collection of 88 instead gives \(N_1=80\) and
\(L_1=-25\), an impairment gain of 25. That negative allowance is a
favorable adjustment to the asset, not a liability owed by the bank.
Paragraph 5.5.14 expressly requires favorable lifetime-loss changes to enter
profit even when expected losses fall below those reflected at acquisition.
Imposing a zero floor would leave carrying amount at 55 and suppress that gain.

The passage of time must also respect the interest basis. Over the next
unchanged year, gross accretion is \(rG\), allowance unwind is \(rL\),
and net interest is their difference:
\[
 rG-rL=r(G-L)=rN.
\]
In the favorable case, these amounts are 5.5, \(-2.5\) and 8.
One balanced decomposition debits the gross asset 5.5 and the signed
allowance account 2.5, and credits interest 8. The allowance becomes
\(-27.5\), so gross 60.5 less allowance \(-27.5\) gives net 88.
The 2.5 is discount unwind within this interest reconciliation; it is not
a second credit-revision gain. The impaired case with net 40 instead earns
interest 4, while gross accretion 5.5 and allowance unwind 1.5 give net 44.

Finally, suppose the favorable case collects only 70 at maturity, with no
reasonable expectation of further recovery. Measuring the remaining net asset
at the known receipt produces a final impairment loss of \(88-70=18\).
Total interest is \(5+8=13\), net impairment gain is \(25-18=7\),
and total profit is \(13+7=20=70-50\). The cash reconciliation agrees:
starting with cash and capital of 200, purchase and collection leave
\(200-50+70=220\), equal to capital plus profit.

Immediately before collection, gross carrying amount is 60.5 and the
signed allowance is \(-9.5\), giving net 70. Receiving cash 70 and
removing those two balances clears the financial asset without another
gain. The contractual amount not collected is nevertheless \(121-70=51\).
The no-further-recovery finding supports financial derecognition under
paragraph 5.4.4; it does not establish that the debtor's legal obligation
has been discharged. The record retains that contractual shortfall and
leaves legal release unassessed.

This terminal-claim example supplies credit forecasts to isolate their
accounting consequences. It does not estimate their probabilities or accuracy.
Intermediate payments, altered contractual terms, fees and earlier write-offs
need their own dated transitions. Positive expected recovery also cannot
justify postponing a required partial write-off. The original credit-adjusted
rate and the distinction between interest
and subsequent credit changes remain the starting conditions for those methods.
